% Section 6 -- Counterfactual policies
% Opening, 6.1 (Table tab:policy-renovation) and 6.2 (Table
% tab:policy-purchase) drafted 5 November 2026 from the outline
% (06-counterfactual-policies-outline.tex, whose opening, 6.1 and 6.2 are
% fenced off with \iffalse there). 6.3 drafted 5 November 2026 (energy
% prices; the paragraph on the price premium waits for its training). The
% close is still the outline, typeset below. The 65--74 band of the policy runs is accepted (5 November
% 2026); no 65--75 re-simulation.
\section{Counterfactual Policies}\label{sec:policy}

Section~\ref{sec:results} shows that older owners live in less efficient
homes because they buy less often. In this section we analyze the effects of
different policies on the efficient share among young and old owners, as well as overall ownership of efficient homes. We focus on three types of policies, distinguishing them by the route through which they act. The first policy focuses on lowering the cost of renovating in place, the fixed cost $\xi_f$ and the
variable cost $\xi$. This policy lowers the threshold $c_j^*$ in
equation~\eqref{eq:cstar-inplace}, the consumption at which an owner who
stays renovates. The second policy acts at a purchase. A point-of-sale
standard requires every home that is bought to be made efficient. Similarly, lower
costs of buying and selling, $\zeta_b$ and $\zeta_s$, make households buy
more often. The last group, higher energy prices, raises the payoff of
efficiency each year without changing either cost. For each policy we
report how the efficient share changes among all owners, among young and
old owners, and in the gap between them.  For the subsidies to
renovation, we also report the public expenditure they need.


We compute the counterfactuals with the approach of
Section~\ref{sec:solution}, in which parameters are inputs of the
networks. Here the inputs are six parameters that a policy can change: the
fixed and variable costs of a renovation, $\xi_f$ and $\xi$; the costs of
buying and selling, $\zeta_b$ and $\zeta_s$; and the maximum loan-to-value
and payment-to-income ratios, $\lambda_L$ and $\lambda_D$. Each ranges over
an interval that contains its value in Tables~\ref{tab:fixed-parameters}
and~\ref{tab:estimates} and every value we consider, and every other
parameter is held at its value in those tables. We call this the
counterfactual solution.  Higher energy prices change two
parameters that are not inputs, the energy bills per unit of size $e^R$
and $e^N$. The point-of-sale standard changes the choices: a buyer can no
longer buy an inefficient home without renovating it at once. Every
simulation draws the same 65,536 households and the same
income shocks, taste shocks, and draws that decide whether a home stays
efficient, so two simulations differ only in the policy. We report the
effect of a policy as the change, in percentage points, from the baseline
of the counterfactual solution.


Before discussing the implications of these policies, we outline three caveats which apply to these experiments. First, the model is partial
equilibrium. Hence, the price premium of efficient homes
does not respond to a policy and a buyer can always buy an efficient home at
the same premium irrespective of how many buyers want one. This matters most for a
policy that raises the demand for efficient homes a great deal, such as
the standard. Second, we compare two steady states, one with the policy
and one without: each household lives its whole life under the policy,
from age 23. We do not compute the transition due to computation constraints. Third, we do not compute welfare. The model has no cost of
energy use beyond the bill a household pays, so household welfare would
leave out any gain from lower emissions, and it would need an assumption
about how each subsidy is financed. We report the efficient share and the
flows behind it.

% \open{Appendix~\ref{app:computation} does not yet have the subsection the
% pointer above promises (the two baselines, the training box, the accuracy
% audits of the counterfactual solution).}
% \open{The cross-check that changes carry over, held for the user: the
% point-of-sale standard imposed without re-solving widens the gap by 4.1
% points on each solution (65--74 band), but halving xi_f changes the gap by
% -0.6 on the counterfactual solution (65--74) and by -1.0 on the estimated
% model (65--75, Section 5.2: 9.62 to 8.59).}
% Decided 5 November 2026: Section 6 reports xi_f halved on the
% counterfactual solution, comparable with xi halved, and the text does not
% mention the larger narrowing in the estimated model.

\subsection{Lowering the cost of renovating in place}\label{sec:policy-renovation}
% Numbers: tab-policy-renovation.csv (tab_policy_renovation.R; the text_*
% rows for the in-place rates and flows); the dose footnote from
% fig-policy-dose.csv and, for the in-place rates by dose
% (renovate_stay_takeup_all: 1.394, 2.396, 4.250, 7.706 percent at xi = 0.08,
% 0.06, 0.04, 0.02), the m7-dose-island-s1 cells; share change per point of
% the rate 2.1-2.2 (xi) and 1.8-2.0 (xi_f); the figure itself is not in the
% paper (decided 5 November 2026); the pound amounts from Section 4.5 (xi_f = 0.159 x
% 49,914 = 7,900; the smallest owned home 20,500, of which 7,900 fixed) and
% outlay x 49,914 (0.574% -> 287, 0.123% -> 61 pounds). Outlay description:
% text_renovations_per100hh_base / text_renovations_per100hh (1.01; 2.69 V,
% 1.55 F) and text_subsidy_per_renovation_units (0.2135 V -> 10,657 pounds;
% 0.0793 F -> 3,958, written 3,950 as half of the rounded 7,900); the outlay
% definition from epc_revised/counterfactual_outcomes.py (outlay(): theta_move
% = 1; survival weights from simulation.survival_weights).

A renovation costs $C(h)=\xi_f+\xi p^N h$
irrespective of whether the owner renovates at a
purchase or in place. A cut in either part i.e., the fixed cost $\xi_f$ or the variable cost $\xi$ lowers the threshold $c_j^*$ in
equation~\eqref{eq:cstar-inplace} and narrows the threshold gap in
equation~\eqref{eq:window} which is the extra cost of renovating in place over
buying an efficient home. We halve each part separately. The variable cost
$\xi$ is reduced from 8 to 4 percent of the value of the home, and the fixed cost
$\xi_f$ from 0.159 to 0.079. The variable cost is the larger part of a
renovation. Renovating the smallest owned home, combining fixed and variable parts, costs about \pounds20,500,
of which \pounds12,600 is the variable cost (Section~\ref{sec:estimates}),
and the variable cost rises with the size of the home. Halving the variable cost saves
\pounds6,300 on the smallest owned home and more on a larger one. Halving
the fixed cost saves \pounds3,950 on every home, irrespective of its size.



\begin{table}[t]
  \centering
  \caption[Lowering the cost of renovating in place]{{Lowering the
  cost of renovating in place}}\label{tab:policy-renovation}
  \begin{threeparttable}
    \small
    \pkgtable{tab-policy-renovation}
    \begin{tablenotes}[flushleft]\footnotesize
      \item {\textit{Notes:} Each column halves one part of the cost
      of a renovation, $C(h)=\xi_f+\xi p^N h$, on the counterfactual
      solution and keeps every other parameter at its value in
      Tables~\ref{tab:fixed-parameters} and~\ref{tab:estimates}. Changes
      from the baseline of the counterfactual solution, from 65,536
      simulated households with the same draws in the baseline and in both
      columns. Shares are A--C owners as a percentage of owners in each age
      group; the gap is the share at 25--34 less the share at 65--74.
      Public outlay: the subsidy paid on all renovations, including those
      that would take place without it. }
    \end{tablenotes}
  \end{threeparttable}
\end{table}


Table~\ref{tab:policy-renovation} shows the effects of halving each cost component. Halving the variable cost raises the efficient share by 6.3 percentage
points among all owners, by 2.5 points at 25--34 and by 8.1 points at
65--74. As old owners are more likely to renovate in place, the gap between young and old owners narrows by 5.6 points. The subsidy roughly
triples the fraction of owners of inefficient homes who renovate in place
each year. This increase occurs at both age groups. But it adds most where renovating in place
supplies most of the inflow into efficient ownership. As older owners
buy less often, it supplies a larger part of theirs: 36 percent at 65--75
against 14 percent at 25--34 in the estimated model
(Section~\ref{sec:results-window}). 

Halving the fixed cost acts in the same way but by less magnitude. This is  because it
saves less on each renovation. It raises the efficient share by 1.7 points
among all owners, by 0.9 at 25--34 and by 1.5 at 65--74, and narrows the
gap by 0.6 points (Table~\ref{tab:policy-renovation}, column~(2)).

To compute the public expenditure with this policy, we assume that the government pays the part of the cost that the
policy removes. That means $\xi p^N h/2$ when the variable cost is halved and
$\xi_f/2$ when the fixed cost is halved, where $h$ is the size of the
renovated home. We count households of every age and tenure, renters
included, with each age weighted by the probability of surviving to it, and
 express it as a percentage of mean working-age gross income. The expenditure is therefore the number of renovations per
household a year times the average subsidy per renovation. Without a
subsidy there are 1.0 renovations per 100 households a year. Halving the
variable cost raises this to 2.7, and the average subsidy is about
\pounds10,700 for the smallest owned home. The expenditure turns out to be 0.57 percent of
mean income per household a year which is about \pounds290. Halving the fixed cost
raises renovations to 1.5 per 100 households a year at \pounds3,950 each,
an outlay of 0.12 percent, which is about \pounds60.\footnote{ In the model, taxes and pensions
depend only on earnings (equations~\eqref{eq:pension} and~\eqref{eq:tax}),
so no other item of the public budget responds to the subsidy.}


\subsection{Acting at a purchase}\label{sec:policy-purchase}
% Numbers: tab-policy-purchase.csv (tab_policy_purchase.R). Column (1): the
% m4 cell mandate (trained with the standard) less island_s1. Column (2): the
% average of T less baseline in m5 (seed 1) and m6 (seed 2), decided 5
% November 2026 (T_s1, T_s2 in the CSV: all 4.92, 5.04; gap 3.77, 0.50;
% moves 10.00, 9.93). text_* rows: ownership (-0.30; the standard imposed
% without retraining, island_s1_mask_mandate, -6.11; T +2.92); owner moves in
% percent (-27.4; T +131); purchases per household in percent (-23.2; T
% +118); the in-place rate of D-G owners in percent (-4.0, not quoted: two
% trainings of the baseline differ by -3.9, and the change by age has mixed
% signs; T -34.4, seeds -36.1 and -32.8);
% purchases of A-C homes under the standard (96.0 percent at 25-34, 88.0 at
% 65-74). Owner moves in the baseline: 7.56 (seed 1), 7.64 (seed 2) per 100
% owners. Section 5.2: tab-flows-accounting (P = 15.0 and 4.5; 40-46
% percent). Outlay of (2): 2.6505 x 49,914 = 1,323 pounds. Training noise:
% island_s2 less island_s1, at most 0.78 points in a share.

The second group of policies acts at a purchase, which, as established in  Section~\ref{sec:results-window}, is where most of the inflow
into efficient ownership happens.
Table~\ref{tab:policy-purchase} reports results for two such policies. The first policy is a
point-of-sale standard. We impose the condition that every home that is bought must be efficient. Hence, a
buyer either buys an efficient home or renovates an inefficient one at
purchase. As this imposition is not determined by a parameter of the networks,  we solve the
model again with the standard in place. Households therefore know of this policy
throughout their lives and plan for it. In the second policy we halve the costs
of buying and selling, $\zeta_b$ from 3.9 to 1.95 percent of the value of
the home and $\zeta_s$ from 5.0 to 2.5 percent. Table~\ref{tab:policy-purchase} reports the results.

\begin{table}[t]
  \centering
  \caption[Acting at a purchase]{{Acting at a
  purchase}}\label{tab:policy-purchase}
  \begin{threeparttable}
    \small
    \pkgtable{tab-policy-purchase}
    \begin{tablenotes}[flushleft]\footnotesize
      \item {\textit{Notes:} Column~(1) requires every home that is
      bought to be A--C; the model is solved again with the standard in
      place, on one training of the networks. Two trainings of the baseline
      of the counterfactual solution differ by at most 0.8 points in each
      share. Column~(2) halves $\zeta_b$, from 0.039 to 0.0195, and
      $\zeta_s$, from 0.050 to 0.025, on the counterfactual solution, and
      reports the average of the effects on two trainings of its networks.
      Every other parameter is at its value in
      Tables~\ref{tab:fixed-parameters} and~\ref{tab:estimates}. Changes
      from the baseline of the same training, from 65,536 simulated
      households with the same draws in the baseline and with the policy.
      Shares and the gap as in Table~\ref{tab:policy-renovation}. Owner
      moves are sales by owners.}
    \end{tablenotes}
  \end{threeparttable}
\end{table}

In column (1) we report the results of the standard policy. The standard raises the efficient share by 13.2 points among all owners. Looking at age groups, it raises the share by
16.3 points at 25--34 and by 9.9 points at 65--74. Consequently, it widens the gap by 6.4
points. In the estimated
model, between 40 and 46 percent of purchases are of efficient homes, and a
few more are renovated at purchase.
As under the standard every purchase adds to the efficient stock, the gain
is largest where households buy most often.  Renovation in place changes very little, since the standard does not
apply to owners who stay. 

It is important to note that the standard  makes buying expensive for a household that would have
bought an inefficient home. This is because  it must pay at least $(1+\zeta_b)\Delta p'$
more, the premium of an efficient home and the buying cost on it. Such a
buyer's consumption is below the threshold $c_b^*$ of
Proposition~\ref{prop:buy}, so the efficient state is worth less to it than
this cost, and therefore some of these households put off moving. Consequently, owner moves
reduce by 2.1 per 100 owners a year, or 27 percent, and purchases fall by 23
percent. From this perspective we observe a trade-off. Each move that does not take place is a  purchase window that
does not open. Hence,  the fall in moves partly offsets the gain which the standard intended to generate. However, overall homeownership barely changes. It falls by 0.3 points as households know
of the standard in advance  and plan for it. When the standard is imposed without warning
instead, so that households make their plans as if there were no standard
and meet it only when they buy, ownership actually falls by 6.1 points. Note that the standard
needs no public expenditure as it is a regulation and its cost falls on buyers.



Halving the costs of buying and selling makes households buy more often.
Purchases rise by 118 percent and owner moves rise by
10.0 per 100 owners a year, from 7.6 (Table~\ref{tab:policy-purchase},
column~(2)).  As households move more often, the efficient share rises by 5.0 points among all owners. The share rises by 6.6
points at 25--34 and by 4.5 at 65--74.\footnote{  The two trainings of the networks
agree on the gain among all owners, 4.9 and 5.0 points, but less on how it
splits by age: the gap widens by 3.8 points on one and by 0.5 on the other.}
 Ownership rises by 2.9 points. However, again, there is a trade-off.
Owners who expect to move sooner renovate in place less, since the
threshold $c_j^*$ in equation~\eqref{eq:cstar-inplace} rises with the
probability of moving. Therefore, the fraction of owners of inefficient homes who
renovate in place each year falls by about a third.

\subsection{Changing the payoff}\label{sec:policy-payoff}
% Numbers: tab-policy-payoff.csv (tab_policy_payoff.R; Table tab:policy-payoff
% and the text_* rows): the m4 cells
% energy_price_x2 and energy_gap_zero, each trained on its own network (seed
% 1 only), less island_s1. Shares 5.22 / 7.12 / 4.31, gap 2.81; in percent
% 39.8 and 36.6; the ratio of the 25-34 to the 65-74 share 1.520 -> 1.556
% (island_s2: 1.588). Bill gap removed: -3.50; -23.3 and -25.5 percent.
% Routes: efficient homes bought (per household) +28.3 percent; the in-place
% rate of D-G owners +20.9 percent; the efficient share of purchases 38.4 ->
% 47.9 (island_s2 36.8); the inflow at a purchase 88.7 percent at 25-34 and
% 69.8 at 65-74 (island_s1). Ownership -2.87; by band -0.46 (within the
% variation between trainings), -2.47, -4.34, -3.47, -3.27, -4.03 (island_s2
% less island_s1 at most 1.40, at 35-44). Median rent paid -8.7 percent;
% median value of an owned home to income at 35-64 -4.9 percent; renters
% 35-64 who own two years later 17.4 -> 15.5 percent (island_s2 17.7);
% owners who rent two years later rise at 25-34 (9.0 -> 9.7), 35-64 and 65+.
% Owner moves +0.75 per 100 owners (island_s2 +0.09); per 100 households
% +0.235, purchases +0.167, so at least 71 percent of the extra moves end in
% a purchase (the script's comment gives the bound). The derivatives in the
% footnote follow from eq:cstar-inplace and eq:cstar-buy with C-bar and ell
% rising one for one with Delta e.
% Numbers (price premium doubled, gamma_p 0.03 -> 0.06): tab-policy-payoff.csv
% (tab_policy_payoff.R), from replication/data/model/batch-2026-11-05/out/
% m9-policy-experiments-seed20260926-n65536 (RTX 4090, 6 November 2026):
% premium_x2 less island_s1 of the same batch (island_s1 and island_s2
% re-simulated on the 4090, not the m4 cells used above; the two island_s1
% differ by at most 0.23 points in a share, gap 6.12 in m4 and 6.35 in m9, so
% the text quotes changes, not levels: the gap closes by 5.64/6.35 = 89
% percent, 92 percent of the m4 gap). Shares 13.605 ->
% 8.870 (-4.73), 25-34 18.056 -> 9.721 (-8.34, -46.2 percent), 65-74
% 11.703 -> 9.005 (-2.70, -23.1 percent), gap 6.353 -> 0.716 (-5.64); ratio
% of the 25-34 to the 65-74 share 1.543 -> 1.080 (island_s2: 1.619).
% 35-44 -9.03, 45-54 -6.37, 55-64 -4.96, 75+ -0.64 (island_s2 less
% island_s1: -0.69). Efficient homes bought per household -43.3 percent
% (25-34 -51.5, 65-74 -32.7); efficient share of purchases 38.4 -> 21.9
% (island_s2 36.6); in-place rate of D-G owners +41.0 percent (25-34 +31.2,
% 65-74 +61.4); renovated at purchase 4.45 -> 6.42 percent of purchases;
% homes bought efficient as a share of the inflow (island_s1) 84.1 percent
% at 25-34, 59.2 at 65-74; inflow per household at 65-74 +4.2 percent.
% Ownership +0.13 (25-34 -0.77; island_s2 less island_s1 -0.25); owner moves
% -0.10 per 100 owners (island_s2 +0.10). C(h)/((1+zeta_b) gamma_p h) at
% gamma_p = 0.06: 1.283 + 2.545/h, 2.09 at h = 3.14 and 1.37 at h = 31.1
% (rows text_cost_ratio_*).

The policies in Sections~\ref{sec:policy-renovation}
and~\ref{sec:policy-purchase} alter the cost associated with achieving the efficient state. However, a
policy can instead change the payoff from owning an efficient home. In the model, an efficient home pays its owner in
three ways: (i) comfort,  (ii) a lower energy bill, and (iii) a higher price when it
is sold.
Section~\ref{sec:results-horizon} varied comfort. In this section  we vary the
other two. The energy
bill is a part of the payoff that a policy, such as a tax on household
energy, can alter. A higher energy price raises every bill and, with them,
the yearly saving $\Delta e=(e^N-e^R)h$. We double energy prices for both
qualities, which doubles $\Delta e$. As energy prices are not a parameter of
the networks, we solve the model again with doubled prices. We
then double the price premium. Table~\ref{tab:policy-payoff} reports both
experiments.

\begin{table}[t]
  \centering
  \caption[Changing the payoff]{{Changing the payoff}}\label{tab:policy-payoff}
  \begin{threeparttable}
    \small
    \pkgtable{tab-policy-payoff}
    \begin{tablenotes}[flushleft]\footnotesize
      \item {\textit{Notes:} Column~(1) doubles the energy bill per
      unit of size of both qualities, $e^R$ and $e^N$. Column~(2) doubles
      the price premium of an efficient home, $\gamma_p$, from 0.03 to 0.06.
      }
    \end{tablenotes}
  \end{threeparttable}
\end{table}

Doubling energy prices raises the efficient share by 5.2 points among all
owners (Table~\ref{tab:policy-payoff}, column~(1)). Looking at age groups, the efficient share increases by 7.1 points at 25--34 and by 4.3 points at 65--74. Hence,  the gap
widens by 2.8 points. In proportion, however, the share rises by approximately
the same amount at both ages,  40 percent at 25--34 and by 37 percent at 65--74.  Like comfort in
Figure~\ref{fig:mechanisms-payoff}, the bill saving sets the level and not the slope of the age profile.



The saving does not enter the two thresholds alike. A buyer who takes the
efficient home pays the lower bill from its first year, and the threshold
$c_b^*$ in equation~\eqref{eq:cstar-buy} falls by the whole saving,
$\Delta e'/\theta'$. An owner who renovates now rather than next year gains
only next year's saving, which it collects only if it survives, stays, and
the home remains efficient. The threshold $c_j^*$ in
equation~\eqref{eq:cstar-inplace} falls by
$\Lambda\tilde\phi_j\mu_{j+1}\Delta e/\theta$.\footnote{The owner pays the
bill of the inefficient home in the year of the work, so $\overline C$ and
$\ell$ both rise one for one with $\Delta e$. Differentiating
equation~\eqref{eq:cstar-inplace} then gives $\partial c_j^*/\partial\Delta
e=-\Lambda\tilde\phi_j\mu_{j+1}/\theta$, and differentiating
equation~\eqref{eq:cstar-buy} gives $\partial c_b^*/\partial\Delta
e'=-1/\theta'$.} A larger saving therefore widens the gap between the two
thresholds in equation~\eqref{eq:window}, and buying an efficient home gains
more than renovating in place. The number of efficient homes bought rises by
28 percent, and the share of purchases that are efficient rises from 38 to
48 percent, whereas the rate of renovation in place rises by 21 percent. At
every age most efficient homes are acquired at a purchase, 89 percent of
the inflow at 25--34 and 70 percent at 65--74. Hence the share rises in
about the same proportion at both ages.

Doubling energy prices also has associated costs. Ownership falls by 2.9
points, and by between 2.5 and 4.3 points in every age group from 35 on.
Owners and renters alike pay the bill on every unit of size, so housing
becomes dearer for both. Renters rent smaller homes, as the median rent
falls by 9 percent, and buyers buy smaller ones, as the median value of an
owned home relative to income falls by 5 percent at 35--64. The fall in
ownership is consistent with the floor on the size of an owned home. An
owned home cannot be smaller than $\underline h^{\mathrm{own}}$, 3.14 times
mean income, whereas a rented home can. A household that now wants less
housing than the smallest owned home therefore rents instead. Fewer renters
buy: the share of renters aged 35--64 who own two years later falls from
17.4 to 15.5 percent, and slightly more owners sell and rent. Owners also
move more often, by 0.8 per 100 owners a year, and at least 70 percent of
the extra moves end in another purchase.

The third part of the payoff is the price premium, which an owner collects
when it sells an efficient home and its heirs collect if it dies. A policy
that makes buyers value the certificate more, for example by making the
rating more visible when a home is sold, raises the premium $\gamma_p$. We
double it, from 3 to 6 percent, and, as with energy prices, solve the model
again. The buyers of efficient homes pay the premium to their sellers, so it
needs no public outlay. Doubling the premium lowers the efficient share by
4.7 points among all owners (column~(2)), by 8.3 points at 25--34 and by
2.7 points at 65--74, and the gap narrows by 5.6 points, almost nine-tenths
of it. In proportion, the share falls by 46 percent at 25--34 and by 23
percent at 65--74. Unlike the bill saving, the premium changes the slope of
the age profile. It barely changes who owns or how often owners move.

The premium enters the two routes to efficiency with opposite signs. A buyer
pays it. Doubling it doubles the cost of buying the efficient state,
$C_b=(1+\zeta_b)\Delta p'$, and so lowers the extra one-off cost of
renovating in place, $C(h)-C_b+\Delta e$, which sets the gap between the two
thresholds in equation~\eqref{eq:window}. Renovating the smallest owned home
now costs 2.1 times the extra outlay of buying it efficient, rather than 4.2
times, and renovating the largest 1.4 times rather than 2.7. An owner who
renovates, by contrast, gets more of its cost back, since the premium is
what a sale returns and what the heirs receive: a larger $\Delta p$ shrinks
both the loss at exit and the loss at death in
equation~\eqref{eq:rule-inplace}. The routes in column~(2) show both
responses. The number of efficient homes bought falls by 43 percent, and the share of
purchases that are of efficient homes falls from 38 to 22 percent. The
fraction of owners of inefficient homes who renovate in place each year
rises by 41 percent, and the share of purchases renovated at purchase rises
from 4.4 to 6.4 percent.

The fall in efficient purchases outweighs the rise in renovation, and most
so among the young, who acquire efficiency mainly by buying it. At 25--34,
homes bought efficient make up 84 percent of the inflow into efficient
homes, against 59 percent at 65--74. At 25--34 the number of efficient homes
bought falls by 52 percent, whereas the in-place rate rises by 31 percent.
Older owners respond more on the other margin: at 65--74 the in-place rate
rises by 61 percent, and the inflow into efficient homes rises by 4 percent.
Their share nevertheless falls, because they reach 65 with fewer efficient
homes. The share at 55--64 falls by 5.0 points, and efficiency lasts about
six years. After 75 the share does not change by more than the variation
between two trainings of the networks.

The experiment bears on a natural conjecture, that the gradient reflects a
small premium. When the market pays little for efficiency, an owner who
expects to leave the home soon, by moving or dying, gets little of the cost
of a renovation back. A larger premium would then bring older owners into
efficient homes and close the gap from below. In the model, a larger premium
does raise renovation, most among older owners, but it lowers their
efficient share, and the gap closes from above, through what the young buy.
The premium therefore acts on the gradient through the purchase window of
Section~\ref{sec:results-window}: the larger the premium, the smaller the
advantage of buying efficiency over renovating, and the young use that
route most. For the same reason, the size of the gradient depends on the
premium measured in the data. Because the premium is fixed in the model,
the experiment compares two markets that price efficiency differently; the
supply of efficient homes for sale does not respond.

% % Section 6 -- Counterfactual policies (outline: paragraphs and what each must
% say; no prose yet). Written 29 October 2026.
%
% Length rule (agreed 29 October): the main text runs to about 40--45 pages
% with about 12 exhibits. Each result goes either to the main text, with full
% intuition, or to the online appendix, as a table and one paragraph with a
% one-sentence pointer from the text. Budget for this section: about 8 pages
% and 1--2 exhibits. \tier{MAIN} and \tier{ONLINE APP.} mark the tier of each
% item.
%
% Numbers are from the draft exhibits of the replication package: the policy
% network ("island"), seed 1, n = 65,536, 30 September batch. None of these
% tables is in the paper yet.
%   levers     output/tables/tab-policy-levers.csv   (m5 factorial, m7 doses)
%   variants   output/tables/tab-policy-variants.csv (m4: mandate, energy
%              prices, bill gap, mortality, no A-C for sale; island vs R0)
%   seeds      output/tables/tab-policy-seeds.csv    (m5 and m6)
%   doses      output/figures/fig-policy-dose.csv    (m7)
% The Shapley split by age and the flows by route and age were computed on 29
% October from the m5 and m6 cells in a scratch script. They are not yet in
% any package script.
% Decided 29 October: the section reports only changes from the policy
% network's own baseline, never its levels, and the network is not retrained.
% Every level in the paper comes from the estimated model (Section 5). So no
% ratio of shares and no statement that needs a level (such as the gradient
% reversing); proportional statements use percent changes within the network.
\providecommand{\tier}[1]{\textsf{\small [#1]}~}

% The opening and 6.1 were drafted as prose in 06-counterfactual-policies.tex
% on 5 November 2026; their outline is kept here but not typeset.
\iffalse
\outpara{Opening: the question and the answer (before the first subsection)}
\begin{itemize}
  \item \tier{MAIN} The question: which policies raise efficiency among older
    owners? Section~\ref{sec:results} implies that the route a policy acts on
    decides whom it reaches. Older owners rarely buy, so a policy that acts
    at a purchase reaches mostly the young. A policy that lowers the cost of
    renovating in place acts on the threshold gap, which every owner faces.
  \item \tier{MAIN} The answer up front. Halving the variable cost of
    renovation raises the A--C share of owners at 65--74 by 8.1 points and at
    25--34 by 2.5, narrowing the gap between them by 5.6 points, more than
    half the 9.6-point gap at the estimates (65--75; to be restated in the
    same band after the re-simulation). A point-of-sale standard
    raises the share the most, by 13.2 points, but widens the gap by 6.4.
  \item \tier{MAIN} Road map: 6.1 to 6.4.
  \item \open{Title: ``Counterfactual Policies'' or one that states the
    question, e.g.\ ``Which Policies Reach Older Owners''.}
\end{itemize}

\subsection{How the counterfactuals are computed}

\outpara{One paragraph in the main text; details in the online appendix}
\begin{itemize}
  \item \tier{MAIN} One network holds the estimated parameters fixed and
    takes the policy parameters as inputs: $\xi_f$, $\xi$, $\zeta_b$,
    $\zeta_s$, $\lambda_L$, $\lambda_D$, each trained over a box around its
    baseline. Each policy is then a simulation in which households
    re-optimize. The point-of-sale standard and the energy-price variants
    change the model, not a lever, so each is trained from the same network.
  \item \tier{MAIN} Every policy uses the same draws as the baseline (common
    random numbers). Effects are differences from that baseline, in
    percentage points.
  \item \tier{MAIN} Scope: (i) partial equilibrium: house prices, the price
    premium of efficient homes, and the supply of efficient homes for sale
    are held fixed, which matters most for a large standard; (ii) each
    policy is in place over the whole life cycle, from 23 on, and the
    transition after a reform is not computed; (iii) no welfare
    calculation.
  \item \tier{MAIN, one sentence} The policy network is a second
    approximation of the same model, and its baseline lies below the
    estimates: 13.6 percent of all owners against 16.2. So we report only
    changes from its own baseline, and every level in the paper comes from
    the estimated model. Changes carry over: the standard imposed directly
    widens the gap by 4.1 points on each network.
  \item \tier{ONLINE APP.} Where the two baselines differ: both solve the
    same problem at the same parameters, and the two policy-network seeds
    agree; they differ mainly in whether a renter who buys picks an A--C
    home, a close call that small approximation errors tip. A table of the
    two baselines (shares, gap, ownership, moves), the cross-check of the
    standard, the per-age accuracy audits, and the training box.
    Appendix~\ref{app:computation} covers the network.
\end{itemize}
\fi

% Drafted as prose in 06-counterfactual-policies.tex on 5 November 2026
% (heading, text and Table tab:policy-renovation moved there). The outline
% below is kept for reference; its ``why the old gain most'' bullet is
% superseded: 72 against 22 D--G owners per 100 households reflects
% ownership by age, not the share among owners (82 and 88 percent of
% owners).
\iffalse
\subsection{Lowering the cost of renovating in place}

\outpara{The two costs and what they act on}
\begin{itemize}
  \item \tier{MAIN} The variable cost $\xi$ is halved from 0.08 to 0.04, and
    the fixed cost $\xi_f$ from 0.159 to 0.079. Either lowers the owner's
    threshold $c_j^*$ and narrows the threshold gap
    (Section~\ref{sec:results-window}).
\end{itemize}

\outpara{The variable cost halved (main table)}
\begin{itemize}
  \item \tier{MAIN} All owners $+6.3$; 25--34 $+2.5$; 65--74 $+8.1$. The gap
    narrows by 5.6 points.
  \item \tier{MAIN} Why the old gain most: renovation in place at 65--74
    rises by 2.0 per 100 households a year, from 0.9, and at 25--34 by 0.5,
    from 0.2. The subsidy acts on owners of D--G homes, and they are
    concentrated among the old: 72 per 100 households at 65--74 against 22
    at 25--34. Buying an efficient home hardly changes, by $-0.1$ at both
    ages.
  \item \tier{MAIN} Outlay: 5.7 ($\times10^3$, model money per
    household-year, inframarginal payments included), or 0.91 per point of
    the all-owner share.
\end{itemize}

\outpara{The fixed cost halved (main table)}
\begin{itemize}
  \item \tier{MAIN} All owners $+1.7$; 25--34 $+0.9$; 65--74 $+1.5$; gap
    $-0.6$. Same direction as the variable cost, smaller. Outlay
    0.73 per point.
  \item \open{Why the fixed cost does so much less than the variable cost
    at a similar outlay per point. Check against the size of each cost
    relative to a typical renovation before writing an intuition.}
\end{itemize}

\outpara{Dose (online appendix, one sentence in the text)}
\begin{itemize}
  \item \tier{ONLINE APP.} Figure fig:policy-dose. Effects grow more than in
    proportion to the cut. A cut of the variable cost to 25 percent raises
    the all-owner share by 13.6 points and narrows the gap by 10.6. Whether
    the gradient would reverse is a statement about levels, so it is not
    made.
\end{itemize}
\fi

% 6.2 drafted in 06-counterfactual-policies.tex (5 November 2026).
\iffalse
\subsection{Acting at a purchase}\label{sec:policy-purchase}

\outpara{A point-of-sale standard (main table)}
\begin{itemize}
  \item \tier{MAIN} Every purchase of a D--G home is renovated at purchase.
    All owners $+13.2$, the largest effect of any policy; 25--34 $+16.3$;
    65--74 $+9.9$. The gap widens by 6.4 points. Ownership falls by 0.3
    points, and owner moves fall by 2.1 per 100
    owners.
  \item \tier{MAIN} Why: the standard acts on the window, and the young
    pass through it most often. The old gain largely because they bought
    under the standard earlier in life. This comes from comparing whole life
    cycles (6.1), so a reform introduced today would reach today's old
    owners less.
  \item \open{Outlay: the standard costs the government nothing, and the
    cost falls on buyers. Leave the outlay column blank or report the
    buyers' renovation spending?}
\end{itemize}

\outpara{Buying and selling costs halved (main table, level only)}
\begin{itemize}
  \item \tier{MAIN} All owners $+4.9$. Owner moves rise by 10.0 per 100
    owners, from 7.6. More owners pass through the window: buying an
    efficient home rises by 2.6 per 100 households a year at 25--34 and by
    3.0 at 65--74, while renovation in place falls slightly.
  \item \tier{MAIN} Outlay 5.46 per point, six times that of the variable
    cost.
  \item \tier{MAIN, one sentence; ONLINE APP.} The split by age is not
    robust: the gap changes by $+3.8$ on seed 1 and $+0.5$ on seed 2, so
    only the level is reported (Table tab:policy-seeds).
\end{itemize}
\fi

% 6.3 drafted in 06-counterfactual-policies.tex (5 November 2026; energy
% prices). The rest of this block (the online-appendix variants, the
% combined cost levers) is kept for the online appendix.
\iffalse
\subsection{Changing the payoff}\label{sec:policy-payoff}

\outpara{What shifts the level without tilting it (one paragraph)}
\begin{itemize}
  \item \tier{MAIN} Energy prices doubled: all owners $+5.2$; 25--34
    $+7.1$; 65--74 $+4.3$. Ownership falls by 2.9 points. In percent of the
    policy network's own baseline, the share rises by about 40 percent at
    25--34 and 37 percent at 65--74. Like comfort in
    Figure~\ref{fig:mechanisms-payoff}, a higher payoff per year raises the
    share at every age in about the same proportion.
  \item \tier{ONLINE APP., one sentence each} No energy-bill gap between
    A--C and D--G homes: $-3.5$. Credit limits: debt-to-income
    0.33 to 0.365, 0.0; loan-to-value 0.90 to 0.85, $-0.2$. Mortality halved
    and no A--C homes for sale on the policy network, as a check on
    Section~\ref{sec:results}.
\end{itemize}

\outpara{Combining the cost levers (online appendix, one sentence in the
text)}
\begin{itemize}
  \item \tier{ONLINE APP.} All three halved: $+12.2$ (seed 2: $+12.7$).
    Shapley shares: variable cost 57 percent, fixed cost 26, buying and
    selling costs 17 (seed 2: 54, 25, 21).
  \item \tier{ONLINE APP.} The levers interact. The fixed cost's share, 3.2
    points, is twice its effect alone, 1.7. The share of buying and selling
    costs, 2.1, is less than half its effect alone, 4.9. Cheaper renovation
    makes the fixed cost bite more and a cheaper move matter less.
  \item \tier{ONLINE APP.} Split of the gap: the variable cost $-5.5$ and
    $-5.1$ on the two seeds, the fixed cost $-1.5$ and $-1.4$; buying and
    selling costs $+4.8$ and $+1.6$, not robust.
\end{itemize}
\fi

% Decided 5 November 2026: the section describes the implications of each
% policy and does not rank them. So no outlay per point (report the total
% outlay), and the close below must be rewritten without ``lowest outlay
% per point'' or other ranking language.
\outpara{Section close (one paragraph)}
\begin{itemize}
  \item \tier{MAIN} Whom a policy reaches depends on the route it acts on.
    Lowering the cost of renovating in place reaches older owners and
    flattens the gradient, at the lowest outlay per point. Acting at a
    purchase raises the level the most but widens the gradient. Raising the
    payoff per year raises the share in about the same proportion at every
    age.
  \item \tier{MAIN} Scope, restated briefly: partial equilibrium, whole life
    cycles, no welfare.
\end{itemize}

\outpara{Exhibits and open decisions (planning note, not text)}
\begin{itemize}
  \item \tier{MAIN} One table, tab:policy, of changes only (no level row;
    the notes point to the levels at the estimates in Section 5), with
    numbered columns built in the R script: (1) all owners, (2) 25--34,
    (3) 65--74, (4) gap, (5) owner moves, (6) outlay per point. Rows:
    \begin{itemize}
      \item renovating in place: the variable cost halved, the fixed cost
        halved;
      \item at a purchase: the point-of-sale standard (column 6 a dash: a
        regulation, no public outlay), buying and selling costs halved
        (columns 1, 5 and 6 only);
      \item the payoff: energy prices doubled.
    \end{itemize}
    It merges \texttt{tab\_policy\_levers.R} and
    \texttt{tab\_policy\_variants.R}.
  \item \open{An optional main figure, which would be new: the effect at
    65--74 against the effect at 25--34 for each policy, with the
    45-degree line, so that which policies tilt the profile is visible at
    a glance.}
  \item \tier{ONLINE APP.} Table tab:policy-seeds, extended with the Shapley
    split by age; credit limits, the bill gap, mortality halved and no A--C
    homes for sale (tab:policy-variants-other, extended); the two baselines
    and the cross-check of the standard; Figure fig:policy-dose.
  \item \open{Experiments not yet run, each needing a training on a pod:
    a grant for renovating in place only; a larger price premium through
    disclosure ($\gamma_p$){\draft, in training since 5 November 2026}.
    None is needed for the argument above.}
  \item \open{Welfare by age (from the value networks, new code) or energy
    saved (from the simulations): a JPE referee is likely to ask for one
    of them.}
\end{itemize}

